\documentclass[11pt]{article}

\usepackage[margin=1in]{geometry}
\usepackage{amsmath,amssymb,amsthm,mathtools}
\usepackage{stmaryrd}
\usepackage{xcolor}
\usepackage{hyperref}
\usepackage{listings}
\usepackage{tikz}
\usetikzlibrary{positioning,arrows.meta,calc,fit,backgrounds}
\usepackage{enumitem}

\hypersetup{
  colorlinks=true,
  linkcolor=blue!50!black,
  citecolor=blue!50!black,
  urlcolor=blue!50!black
}

% --- Theorem environments ---
\theoremstyle{plain}
\newtheorem{theorem}{Theorem}[section]
\newtheorem{lemma}[theorem]{Lemma}
\newtheorem{proposition}[theorem]{Proposition}
\newtheorem{corollary}[theorem]{Corollary}

\theoremstyle{definition}
\newtheorem{definition}[theorem]{Definition}
\newtheorem{example}[theorem]{Example}
\newtheorem{construction}[theorem]{Construction}
\newtheorem{algorithm}[theorem]{Algorithm}

\theoremstyle{remark}
\newtheorem{remark}[theorem]{Remark}

% --- Notation ---
\newcommand{\sem}[1]{\llbracket\,#1\,\rrbracket}
\newcommand{\reflect}[1]{\ulcorner #1 \urcorner}
\newcommand{\rew}{\leadsto}
\newcommand{\Names}{\mathcal{N}}
\newcommand{\Procs}{\mathcal{P}}
\newcommand{\Terms}{\mathbb{T}}
\newcommand{\Pos}{\mathbb{P}}
\newcommand{\States}{\mathcal{S}}
\newcommand{\Lhs}{\mathcal{L}}
\newcommand{\hd}{\mathsf{hd}}
\newcommand{\arity}[1]{\##1}
\newcommand{\dom}{\mathcal{D}}
\newcommand{\gcp}{\mathsf{gcp}}
\newcommand{\reduce}{\mathsf{reduce}}
\newcommand{\deriv}{\mathsf{deriv}}
\newcommand{\lift}{\mathsf{lift}}
\newcommand{\nil}{\mathbf{0}}

% Short rho-syntax
\newcommand{\Send}[2]{#1 \,!\, (#2)}
\newcommand{\Recv}[3]{\mathsf{for}\,(#1 \leftarrow #2)\,\{\,#3\,\}}
\newcommand{\Bang}[3]{\mathsf{for}\,(#1 \Leftarrow #2)\,\{\,#3\,\}}
\newcommand{\Let}[3]{\mathsf{let}\;#1 = #2\;\mathsf{in}\;#3}
\newcommand{\New}[2]{\mathsf{new}\;#1\;\mathsf{in}\;\{\,#2\,\}}

\lstdefinelanguage{rho}{
  morekeywords={new,in,for,match,contract,select,if,else,let,Nil},
  sensitive=true,
  morecomment=[l]{//},
  morestring=[b]"
}
\lstset{
  language=rho,
  basicstyle=\ttfamily\small,
  keywordstyle=\bfseries\color{blue!50!black},
  commentstyle=\itshape\color{gray},
  stringstyle=\color{red!60!black},
  breaklines=true,
  showstringspaces=false,
  xleftmargin=1em
}

\title{From Graph-Structured Lambda Theories to Rholang:\\
       Compiling Set-Automaton Interpreters\\ with Inherent Concurrency}
\author{}
\date{\today}

\begin{document}
\maketitle

\begin{abstract}
A finite presentation of a graph-structured lambda theory (GSLT) specifies a
small-step operational semantics by a finite signature of term constructors,
a finite set of equations, and a finite inductive description of a reduction
relation. We give a complete algorithm that, from such a presentation,
produces a Rholang program implementing the corresponding small-step
interpreter. The interpreter is structured as a Bouwman--Erkens--Groote
\emph{set automaton} for pattern matching against the left-hand sides of the
presentation, paired with a fixed schedule of substitution and rewrite-rule
firing. Each state of the set automaton is reflected to a Rholang channel;
each transition is realised by a persistent \textsf{for}-receive on that
channel; the indeterminate order of processing position pointers is realised
by the indeterminate order of message delivery on those channels. We prove a
\emph{parallelism preservation} theorem: the asynchronous, name-passing
semantics of Rholang makes every parallelism opportunity exposed by the set
automaton --- both intra-term (independent position pointers) and inter-term
(independent redex firings) --- directly available to the Rholang runtime
without further scheduling glue. The compiler is total on any finitely
presented GSLT, the produced program is of size polynomial in the automaton,
and the resulting interpreter inspects each function symbol of the subject
term exactly once per redex search.
\end{abstract}

\tableofcontents

\section{Introduction}\label{sec:intro}

Operational semantics for languages built on term rewriting typically
consists of a finite signature, a finite set of equations, and a finite
inductive specification of a reduction relation. The \emph{interpreter}
that implements such a semantics has two well-understood components:
\begin{enumerate}[leftmargin=*]
  \item a \emph{matcher} that finds all redex occurrences in the subject
        term --- the positions in the term where a left-hand side of some
        rewrite rule matches; and
  \item a \emph{rewriter} that fires a chosen redex by substituting the
        instance of the corresponding right-hand side and announces the
        positions in the resulting term where matching should resume.
\end{enumerate}
Both components are eminently parallelisable: distinct positions in a term
can be examined independently, and distinct redex sites in a non-overlapping
rule set can be fired independently.

Two contributions of the last decade have sharpened the matcher side of this
picture. First, the \emph{set automaton} of Erkens and
Groote~\cite{ErkensGroote2021} solves the subterm pattern-matching problem
for a finite pattern set against a finite signature in a way that visits
each function symbol of the subject term at most once and exposes the
algorithm's parallel structure intrinsically: at each state, the automaton
emits a set of independent (state, position)-pairs that may be processed in
any order, or simultaneously, by distinct workers. Second, the rewriter
side has been put on a categorical footing by \emph{graph-structured lambda
theories} (GSLTs), which view an operational semantics as a category with
finite limits equipped with a distinguished object of processes and a
distinguished subobject for the reduction relation, presented finitely by
generators, equations and inductive clauses.

This paper joins these two threads through Rholang, a concurrent
process-calculus language~\cite{MeredithRadestock2005} whose primitives are
asynchronous send, persistent name-restricted receive, and reflective
name/process quoting. Concretely, we give an algorithm
\[
  \mathsf{compile} \;:\; \mathsf{GSLTPresentation}
                  \;\longrightarrow\; \mathsf{RholangProgram}
\]
with the following properties.
\begin{description}[leftmargin=2em,style=nextline]
  \item[Totality.] $\mathsf{compile}$ is defined on every finitely presented
       GSLT, terminates, and runs in time polynomial in the size of the
       generated set automaton.
  \item[Correctness.] The Rholang program produced by $\mathsf{compile}$
       implements the small-step semantics of the input presentation: for
       every input term encoded as a Rholang process, the program announces
       on a public channel exactly the set of redex (rule, position)-pairs
       given by the presentation, fires them in some sound order, and
       continues until normal form.
  \item[Parallelism preservation.] Every parallel opportunity exposed by the
       set automaton --- independent successor positions, independent
       redexes, independent rewrites --- is materialised as a parallel
       composition or as a free choice of message order in the Rholang
       output, without scheduling glue. The runtime is responsible for
       choosing an interleaving, and any choice is sound.
\end{description}

The technical core is the \emph{reflection of automaton states as channels}.
A set-automaton state is, by construction, a finite set of match obligations.
We reflect this state, as a piece of syntactic data, to a Rholang name; the
transitions out of the state become persistent receives on that name, the
emission of (state, position)-successors becomes a parallel composition of
sends to the successor states' names, and the announcement of a pattern
match becomes a send on a public matches channel. The algorithm thus has a
particularly compact closed form: every artefact of the compiler is one of a
handful of \emph{macros}, each templated by a state or a function symbol of
the signature.

\paragraph{Outline.}
Section~\ref{sec:prelim} fixes notation and recalls the three ingredients:
GSLTs and their finite presentations (\S\ref{ssec:gslt}), the set automaton
of Erkens--Groote (\S\ref{ssec:set-automaton}), and the fragment of Rholang
we use (\S\ref{ssec:rholang}). Section~\ref{sec:encoding} describes the
encoding of subject terms as Rholang processes and of automaton states as
Rholang names. Section~\ref{sec:compile} gives the compilation algorithm
proper, organised as nine code-emission rules. Section~\ref{sec:parallel}
states and proves the parallelism-preservation theorem.
Section~\ref{sec:example} works through a small, complete example.
Section~\ref{sec:discussion} discusses garbage collection, dynamic rule
sets, and other practical concerns. Section~\ref{sec:conclusion} concludes.

\section{Preliminaries}\label{sec:prelim}

\subsection{Graph-structured lambda theories}\label{ssec:gslt}

A graph-structured lambda theory is a categorical packaging of an operational
semantics: the category is the host of term constructors, the
distinguished object is the universe of processes, and the distinguished
subobject is the reduction relation.

\begin{definition}[Graph-structured lambda theory]\label{def:gslt}
A \emph{graph-structured lambda theory} (GSLT) is a tuple
$\mathcal{G} = (\mathcal{C}, P, R)$ where
\begin{itemize}[leftmargin=*]
  \item $\mathcal{C}$ is a category with finite limits;
  \item $\mathcal{C}$ is closed with respect to the cartesian product
        (so internal homs $A \Rightarrow ({-})$ exist as right adjoints);
  \item $P \in \mathcal{C}$ is a distinguished object, the
        \emph{object of processes};
  \item $R \hookrightarrow P \times P$ is a distinguished subobject,
        the \emph{reduction relation}.
\end{itemize}
\end{definition}

For implementation purposes we restrict to GSLTs given finitely.

\begin{definition}[Finite presentation]\label{def:presentation}
A \emph{finite presentation} of a GSLT consists of
\begin{enumerate}[leftmargin=*]
  \item a finite set $\mathcal{S}$ of \emph{base sorts} (in the sense
        of a multisorted Lawvere theory~\cite{TrimbleNLabLawvere}); we
        write
        $\mathcal{T}(\mathcal{S})$ for the closure of $\mathcal{S}$
        under finite products $\times$ and exponentials $\Rightarrow$,
        and call its elements \emph{types};
  \item a finite ranked alphabet
        $\mathbb{F} = \mathbb{F}_0 \cup \mathbb{F}_1 \cup \cdots \cup \mathbb{F}_n$
        of \emph{function symbols} (term constructors), each with a
        fixed arity, a fixed product-shaped input/output profile
        $f : A_1 \times \cdots \times A_{\arity{f}} \to B$ with
        $A_1, \ldots, A_{\arity{f}}, B \in \mathcal{T}(\mathcal{S})$
        (in particular each $A_i$ may itself be a product or an
        exponential --- for example $\mathrm{lam} : (P \Rightarrow P) \to P$
        for the $\lambda$-abstraction constructor, with input type
        the exponential $P \Rightarrow P$), and a finite set
        $\lambda_f$ of \emph{well-formedness equations}
        $g_i(\vec{x}) = h_i(\vec{x})$ between $\mathbb{F}$-terms over the
        inputs $\vec{x} = (x_1, \ldots, x_{\arity{f}})$, which the
        inputs must satisfy modulo $E$ for $f(\vec{x})$ to be
        well-formed;\footnote{Definition~\ref{def:gslt} asks for $\mathcal{C}$
        to have all finite limits, so in principle a function-symbol
        profile could be any finite limit $\mathrm{lim}\,D_f \to B$, not
        merely a product. By the standard product--equaliser
        factorisation of finite limits, every such profile is
        equivalent data to a product profile $f^{\flat} : A_1 \times
        \cdots \times A_n \to B$ together with the finite set of
        equations $\lambda_f$ defining the equaliser. We adopt that
        factorised form for compilation; no expressive power is lost.
        When $\lambda_f = \emptyset$ the profile is a plain product, as
        in a standard ranked alphabet.};
  \item a finite set $E$ of equations between $\mathbb{F}$-terms,
        together with a reduction order $\succ$ on terms that orients
        $E$ into a convergent (confluent and terminating) rewrite system
        $E^{\downarrow}$ via Knuth--Bendix completion, giving canonical
        representatives of equivalence classes; we further require that
        $E^{\downarrow}$ has the \emph{finite variant property} (FVP) of
        Escobar--Meseguer--Sasse~\cite{EscobarMeseguerSasse2008}: for every
        $\mathbb{F}$-term $t$ the set
        \[
          \mathsf{Var}_{E}(t)
          \;:=\;
          \bigl\{ (t\sigma)\!\downarrow_{E^{\downarrow}}
                  \;:\; \sigma \text{ a substitution} \bigr\}
        \]
        is finite up to renaming of metavariables. FVP makes
        $E$-matching --- ``does $P\sigma =_E T$ for some $\sigma$?'' ---
        decidable via folding variant
        narrowing~\cite{EscobarSasseMeseguer2012}, and is enjoyed by all
        standard structural theories (AC, ACI, abelian groups,
        exclusive-or);\footnote{Without FVP, $E$-matching is in general
        undecidable even when $E^{\downarrow}$ is convergent, so KB
        alone is insufficient for the compiler to decide whether a
        subject term has an $E$-equivalent matching a rule's left-hand
        side. Alternative sufficient conditions (left-linearity plus
        termination of basic narrowing; shallow $E$; right-ground $E$)
        also work, but FVP is the cleanest single hypothesis.}
  \item a finite set $\mathcal{R}_{\mathrm{rule}}$ of \emph{rewrite rules}
        $L_j \rew R_j$, where $L_j, R_j$ are morphisms in the syntactic
        Lawvere theory of $\mathbb{F}$ with a common typed domain
        $A_1 \times \cdots \times A_m$ (the \emph{metavariable inputs}
        of the rule) and codomain $P$, and a finite set
        $\mathcal{R}_{\mathrm{ctx}}$ of \emph{contextual rewrites} of the
        shape
        \[
          L_1 \rew R_1, \ldots, L_k \rew R_k
            \;\Rightarrow\; K[L_1,\ldots,L_k] \rew K'[R_1,\ldots,R_k],
        \]
        with $K, K'$ $k$-hole $\mathbb{F}$-morphisms and each $L_i, R_i$
        an $\mathbb{F}$-morphism in the same sense.

        We display rules in the standard tree form, writing repeated
        metavariable tokens for the same input. A repeated occurrence
        of a metavariable $\omega : A$ on the LHS denotes composition
        with the diagonal $\Delta_{A} : A \to A \times A$ (or, more
        generally, an iterated $\Delta$ for several occurrences); the
        LHS as a whole is the categorically linear morphism that, after
        $\Delta$-expansion, uses each metavariable input once. This
        unifies the two traditions of ``linearity'':
        \emph{categorical linearity} (the LHS is a well-formed morphism
        in the syntactic category, with $\Delta$ available because of
        finite products) is always assumed, and \emph{syntactic
        linearity} (no metavariable token repeats) is the special case
        of categorical linearity in which no $\Delta$-composition
        occurs;
  \item a \emph{reduction relation} $R \hookrightarrow P \times P$,
        specified as the least subobject closed under the following
        clauses, read pointwise on ground terms:
        \begin{description}[leftmargin=2em,style=nextline]
          \item[(Rule)] for every $(L_j \rew R_j) \in \mathcal{R}_{\mathrm{rule}}$
               and every ground substitution $\sigma$,
               $(L_j^{\sigma}, R_j^{\sigma}) \in R$;
          \item[(Context)] for every contextual rewrite
               \[
                 L_1 \rew R_1, \ldots, L_k \rew R_k
                   \;\Rightarrow\; K[L_1,\ldots,L_k] \rew K'[R_1,\ldots,R_k]
               \]
               in $\mathcal{R}_{\mathrm{ctx}}$ and every ground substitution
               $\sigma$, if $(L_i^{\sigma}, R_i^{\sigma}) \in R$ for all
               $1 \le i \le k$, then
               \[
                 \bigl(K^{\sigma}[L_1^{\sigma},\ldots,L_k^{\sigma}],\;
                       K'^{\sigma}[R_1^{\sigma},\ldots,R_k^{\sigma}]\bigr) \in R.
               \]
        \end{description}
\end{enumerate}
We write
$\mathcal{R} := \mathcal{R}_{\mathrm{rule}} \cup \mathcal{R}_{\mathrm{ctx}}$
for the full rule set of the presentation.
\end{definition}

The presentation generates a GSLT freely by taking $\mathcal{C}$ to be the
syntactic category of the simply-typed $\lambda$-theory with constants
$\mathbb{F}$, modulo $E$ and $\beta\eta$; $P$ to be a base sort, and $R$ to
be the least subobject closed under (Rule) and (Context).

\paragraph{Binders.}
A function symbol $f \in \mathbb{F}_n$ may carry a set
$\beta_f \subseteq \{1, \ldots, n\}$ of \emph{binding positions}. Under the
internal hom of Definition~\ref{def:gslt}, the $i$-th argument of $f$ for
$i \in \beta_f$ has type $P \Rightarrow P$. Operationally, the subterm at
position $p.i$ of $f(t_1, \ldots, t_n)$ for $i \in \beta_f$ is the opened
body $t_i(z_{p.i})$ for a fresh free name $z_{p.i}$; the position-domain,
sub-pattern, and head-symbol functions extend by the obvious recursion
through opened bodies. Free names introduced by binders behave as leaves in
the term tree, analogous to variables. The set-automaton machinery and the
compiler we describe are agnostic to binders: positions and head symbols
are the only operationally relevant features, and binders contribute only
to the way ground substitutions interact with rule firings (cf.\ \S\ref{ssec:rewriter}).

\paragraph{Patterns, left-hand sides, and sharing constraints.}
A \emph{raw pattern} is a tree in $\Terms(\mathbb{F}, \mathcal{V})$ over a
countable set $\mathcal{V}$ of metavariables, with metavariables permitted
to repeat. The \emph{linearisation} of a raw pattern $\ell$ is the pair
$(\ell^{\flat}, \kappa_{\ell})$ where $\ell^{\flat}$ replaces every
occurrence of every metavariable by a distinct fresh variable, and
$\kappa_{\ell}$ is the set of position-pairs $(p_1, p_2)$ such that the
$p_1$-leaf and $p_2$-leaf of $\ell$ carry the same original
metavariable. We call $\kappa_{\ell}$ the \emph{sharing constraint} of
$\ell$; it is empty iff $\ell$ is syntactically linear. The set of
left-hand sides of a presentation is
\[
  \Lhs \;=\; \{(\ell^{\flat}, \kappa_{\ell}) : \ell \in \mathit{LHSes}(\mathcal{R}_{\mathrm{rule}}) \cup \mathit{LHSes}(\mathcal{R}_{\mathrm{ctx}})\},
\]
and $\Lhs$ is the input to the set-automaton construction below. We
typically drop the $\flat$ superscript and write $\ell$ when the
distinction is clear from context.

\subsection{The set automaton}\label{ssec:set-automaton}

The set automaton of Erkens and Groote~\cite{ErkensGroote2021} solves the
\emph{subterm pattern-matching problem}: given a finite pattern set $\Lhs$,
locate every $(\ell, p)$-pair such that pattern $\ell \in \Lhs$ matches the
subject term at position $p$. We recall the construction at the level of
detail needed for compilation; for proofs and intuitions we refer to the
original paper.

\paragraph{Match obligations and goals.}
A \emph{match obligation} on $\Lhs$ is a finite, non-empty subset of
$\mathsf{sub}(\Lhs) \times \Pos$, where $\mathsf{sub}(\Lhs)$ is the set of
sub-patterns of patterns in $\Lhs$ excluding the metavariable $\omega$, and
$\Pos$ is the set of positions (finite sequences of positive integers, with
$\epsilon$ for the root and $p.q$ for concatenation). A \emph{match
announcement} is a pair $(\ell, p) \in \Lhs \times \Pos$. A \emph{match
goal} is an arrow $mo \to ma$ from a match obligation to a match
announcement, written
\[
  \ell_1 @ p_1, \ldots, \ell_n @ p_n \;\to\; \ell @ p,
\]
read: \emph{in order to announce that $\ell$ matches at $p$, we must
observe sub-pattern $\ell_i$ at $p_i$ for each $1 \le i \le n$.}

\paragraph{The tuple.}
A \emph{set automaton} for $\Lhs$ is a tuple
$M = (\States, s_0, L, \delta, \eta)$ where:
\begin{itemize}[leftmargin=*]
  \item $\States$ is a finite set of states, each state itself a finite
        non-empty set of match goals;
  \item $s_0 = \{\ell @ \epsilon \to \ell @ \epsilon \mid \ell \in \Lhs\}$
        is the initial state;
  \item $L : \States \to \Pos$ is the \emph{state-label function},
        assigning to every state a position in some match obligation
        (the position to inspect next);
  \item $\delta : \States \times \mathbb{F} \to \mathcal{P}(\States \times \Pos)$
        is the \emph{transition function};
  \item $\eta : \States \times \mathbb{F} \to \mathcal{P}(\Lhs \times \Pos \times \mathcal{K})$
        is the \emph{output function}, emitting triples $(\ell, p, \kappa)$
        where $\kappa \in \mathcal{K}$ is the sharing constraint inherited
        from $\ell$, instantiated at the announcement position $p$.
\end{itemize}
The empty set is the \emph{final state}: no transitions, no output.

\paragraph{Operation.}
The automaton is run with a \emph{position pointer} $p$ ranging over
positions in the subject term $t$. A configuration $(s, p)$ at state $s$
and pointer $p$ proceeds by reading the head symbol
$f := \hd(t[p.L(s)])$ at position $p.L(s)$ in $t$; then:
\begin{itemize}[leftmargin=*]
  \item for every $(\ell, p', \kappa) \in \eta(s, f)$, announce that
        pattern $\ell$ matches at position $p.p'$, conditioned on the
        sharing constraint $\kappa$ holding at $p.p'$ in $t$; if $\kappa$
        is empty the announcement is unconditional;
  \item for every $(s', p'') \in \delta(s, f)$, recursively process
        configuration $(s', p.p'')$.
\end{itemize}
The order in which the successor configurations
$\{(s', p.p'') : (s', p'') \in \delta(s, f)\}$ are processed is unspecified.
Sequential implementations use a stack or queue, giving DFS or BFS; the
crucial observation~\cite{ErkensGroote2021} is that any order is sound and
visits each function symbol of $t$ at most once.

\paragraph{Construction.}
The states, label function, and transition function are jointly defined as
a least fixed point. Given a state $s$ and a function symbol
$f \in \mathbb{F}$, the \emph{$f$-derivative} of $s$ is the next-step state
obtained by updating all match obligations in $s$ that mention $L(s)$ under
the observation of $f$ at $L(s)$:
\[
  \deriv(s, f) \;=\; \mathit{unchanged} \;\cup\; \mathit{reduced} \;\cup\; \mathit{fresh},
\]
where (using the obvious component-wise reduction operator on obligations)
\begin{align*}
  \mathit{unchanged} &= \{mo \to ma \in s \mid L(s) \notin \mathsf{pos}(mo)\},\\
  \mathit{reduced}   &= \{\reduce(mo, f, L(s)) \to ma \mid mo \to ma \in s,\;
                          \exists\, (\ell @ L(s)) \in mo : \hd(\ell) = f,\;
                          \reduce(mo, f, L(s)) \ne \emptyset\},\\
  \mathit{fresh}     &= \{\ell @ L(s).i \to \ell @ L(s).i \mid \ell \in \Lhs,\; 1 \le i \le \arity{f}\}.
\end{align*}
The derivative is then partitioned into dependency-equivalence classes (two
obligations are dependent iff their position-sets intersect, and the
equivalence is the transitive closure of dependency); each class becomes a
new state after \emph{lifting}, i.e.\ replacing each position by its image
under the greatest-common-prefix offset. Lifting keeps the position
pointers bounded and the state set finite. The construction is
deterministic up to the (immaterial) choice of $L(s)$ when several root
obligations are available.

\paragraph{Sharing constraints.}
The construction above is unchanged from the syntactically linear case
of Erkens--Groote~\cite{ErkensGroote2021}: it runs on $\Lhs$ viewed as
a set of linearised raw patterns $\ell^{\flat}$, ignoring the sharing
constraints $\kappa_{\ell}$. The constraints reappear at output time:
when state $s$ has a root goal whose announcement fires under symbol
$f$ at position $p$, the triple emitted is $(\ell, p', \kappa)$ where
$\kappa$ is $\kappa_{\ell}$ with every position $q$ in its pairs
replaced by $p.q$ (i.e., reanchored to the announcement position in the
subject term). The set automaton is therefore agnostic to $\Delta$: it
detects syntactic shape and defers structural-equality checking to the
fire stage (\S\ref{ssec:rewriter}).

\paragraph{Example: the rho COMM rule.}
The rho-calculus communication rule
\[
  \mathsf{for}(y \!\leftarrow\! x)\{B\} \mid x!(Q)
    \;\rew\; B[@Q / y]
\]
is the prototypical non-syntactically-linear LHS: as a morphism in the
syntactic Lawvere theory it has shape $N \times (N \Rightarrow P) \times P \to P$
with one channel-input $x$, one body-input $B$ (binding $y$), and one
payload $Q$, but its tree display has two leaves labelled $x$
(composition with $\Delta_N$). The linearisation introduces fresh
metavariables $x_1, x_2 : N$ at positions $1.1$ (subject of the
receive) and $2.1$ (subject of the send), giving sharing constraint
\[
  \kappa_{\mathrm{COMM}} \;=\; \{(1.1,\, 2.1)\}.
\]
The set automaton is built on $\ell^{\flat}_{\mathrm{COMM}}$ in the
usual way; at announcement of $\ell^{\flat}_{\mathrm{COMM}}$ at some
position $p$ in the subject, the emitted triple is
$(\mathrm{COMM}, p, \{(p.1.1,\, p.2.1)\})$, and the fire contract
(\S\ref{ssec:rewriter}) verifies $p.1.1 \equiv p.2.1$ in the subject
before firing.

\subsection{Rholang}\label{ssec:rholang}

We use a small fragment of Rholang based on the rho calculus of Meredith
and Radestock~\cite{MeredithRadestock2005}. Processes $P$ and names $x$
are mutually defined by the grammar
\[
  P ::= \nil \mid \Send{x}{P} \mid \Recv{y}{x}{P} \mid \Bang{y}{x}{P}
        \mid P \mid P \mid {*x} \mid \New{x}{P},
  \qquad
  x ::= @P,
\]
with the following intended readings:
\begin{itemize}[leftmargin=*]
  \item $\nil$: the inert process;
  \item $\Send{x}{P}$: asynchronous send of process $P$ on channel $x$;
  \item $\Recv{y}{x}{P}$: \emph{linear} receive on channel $x$, binding
        the incoming name to $y$ in $P$ and then consuming the receiver;
  \item $\Bang{y}{x}{P}$: \emph{persistent} receive (\textsf{contract}),
        which fires on every send to $x$, binding $y$ in $P$, and remains
        available;
  \item $P \mid Q$: parallel composition;
  \item $*x$: dereference of name $x$, evaluating to the process it
        names;
  \item $\New{x}{P}$: name restriction, introducing a fresh name $x$ in
        $P$;
  \item $@P$: the \emph{reflection} of process $P$ to a name.
\end{itemize}
Structural congruence $\equiv$ contains $\alpha$-equivalence, associativity
and commutativity of $\mid$ with $\nil$ as unit, and the reflection law
$*@P \equiv P$. Names are quotients of processes under $\equiv$. The
principal reduction is
\[
  \Recv{y}{x}{P} \mid \Send{x'}{Q} \;\rew\; P[@Q / y]
  \quad\text{whenever}\quad {*x} \equiv {*x'},
\]
with the persistent variant leaving a fresh copy of the receive in place.

We use Rholang's standard surface syntax in displayed code, writing
\texttt{x!(P)} for $\Send{x}{P}$, \texttt{for(y <- x) \{P\}} for the linear
receive and \texttt{for(y <= x) \{P\}} for the persistent variant, and
\texttt{new x in \{P\}} for $\New{x}{P}$.

\paragraph{Match.}
Surface Rholang adds a \texttt{match} construct on quoted-process patterns.
We use \texttt{match} only to deconstruct the surface symbol and the
arguments of an encoded subject term (\S\ref{sec:encoding}); it is
unnecessary for the set-automaton dispatch, which is realised entirely by
channel names.

\paragraph{Reflection bracket.}
Throughout the paper $\reflect{\cdot}$ denotes a fixed section of $@$
chosen so that distinguishable mathematical objects map to distinct
names: $\reflect{a} =_{\Names} \reflect{b}$ iff $a$ and $b$ are equal as
mathematical objects. We use the bracket on automaton states, term
positions, and signature symbols, generating canonical and distinct
channel names from each.

\section{Encoding terms and states}\label{sec:encoding}

The compiler emits a fixed set of definitions independent of the input
term, plus an encoding of each subject term as a Rholang process. We
describe the encodings here; the compilation algorithm proper is in
\S\ref{sec:compile}.

\subsection{Subject terms}\label{ssec:term-encoding}

Let $t \in \Terms(\mathbb{F})$ be a closed subject term. We encode $t$ as a
Rholang process on a \emph{term channel} $\tau$ that, when read, yields a
constructor tag together with a tuple of names for the immediate
subterms. Concretely, for each function symbol $f \in \mathbb{F}_n$ we emit
a persistent receive that responds to a request on $\tau$ with the head
symbol $f$ and a tuple of names $(\tau_1, \ldots, \tau_n)$ whose targets
are the recursively encoded subterms.

\begin{construction}[Term encoding]\label{cons:term-encoding}
The encoding $\sem{t}_{\tau}$ of a term $t$ on a channel name $\tau$ is
defined inductively:
\begin{itemize}[leftmargin=*]
  \item For $a \in \mathbb{F}_0$ (a constant):
  \[
    \sem{a}_{\tau} \;=\; \Bang{r}{\tau}{\Send{r}{\reflect{a}}}.
  \]
  \item For $f \in \mathbb{F}_n$ with $n \ge 1$ and subterms $t_1, \ldots, t_n$:
  \[
    \sem{f(t_1, \ldots, t_n)}_{\tau}
    \;=\;
    \New{\tau_1, \ldots, \tau_n}{
      \sem{t_1}_{\tau_1} \mid \cdots \mid \sem{t_n}_{\tau_n}
      \mid \Bang{r}{\tau}{\Send{r}{(\reflect{f}, \tau_1, \ldots, \tau_n)}}
    }.
  \]
\end{itemize}
\end{construction}

A query on $\tau$ thus returns the \emph{shape} of $t$ at the root --- its
head symbol plus channel names addressing its arguments --- and leaves
sub-shapes available behind those argument channels. This makes the
encoding lazy in a useful sense: the automaton consults only the positions
it actually needs to inspect.

\paragraph{Binders.}
A function symbol $f \in \mathbb{F}_n$ with binding positions
$\beta_f \subseteq \{1, \ldots, n\}$ encodes binding arguments as
$(P \Rightarrow P)$-valued, i.e.\ as Rholang names that, when sent a fresh
\emph{bound-name token}, respond with the encoding of the opened body
under that token. Concretely, position $i \in \beta_f$ of $f$ is encoded
on a channel $\tau_i$ as a persistent receive
\[
  \Bang{(z, r)}{\tau_i}{\Send{r}{\sem{t_i(z)}_{r}}},
\]
which takes a fresh free name $z$ and a response channel $r$, and emits
the encoded opened body. The automaton invokes this protocol only at
positions in $\beta_f$ that it actually inspects.

\subsection{Automaton states as channels}\label{ssec:state-channels}

For each state $s \in \States$ of the set automaton built from $\Lhs$, the
compiler reserves a single \emph{state channel} $\sigma_s := \reflect{s}$.
The state channel carries messages of the form $(\tau, p)$, where $\tau$
is the channel encoding the subject term and $p$ is a position pointer
into $t$.

A persistent receive on $\sigma_s$ implements the per-state behaviour:
the position $L(s)$ is added to $p$ to obtain the position to inspect;
the term is queried at that position, yielding a head symbol $f$ and the
argument channels; $\eta(s, f)$ is announced on a public \emph{matches}
channel; and each $(s', p'') \in \delta(s, f)$ is sent forward as
$(\tau, p.p'')$ on $\sigma_{s'}$. Section~\ref{sec:compile} renders this
behaviour as explicit Rholang code.

\subsection{Position pointers}\label{ssec:pointer-encoding}

A position pointer $p \in \Pos$ is encoded as a Rholang list of positive
integers, or, equivalently, as a name $\reflect{p}$. Concatenation
$p.q$ is a Rholang list operation. Lookup at a position is realised by a
sequence of \emph{traversal} sends down the argument channels: from a
term channel $\tau$ and a pointer $p = i_1 . i_2 . \cdots . i_k$, the
sequence
\[
  \Send{\tau}{r_0},\;
  \mathsf{for}\, ((f_0, \vec{\tau}_0) \leftarrow r_0)\,
  \{\, \Send{\vec{\tau}_0[i_1]}{r_1},\,
       \mathsf{for}\,((f_1, \vec{\tau}_1) \leftarrow r_1)\, \cdots \,\}
\]
descends the term tree along $p$. This traversal is generic and is emitted
once by the compiler as a contract (the \emph{traversal contract}, see
\S\ref{ssec:emit-traversal}).

\section{The compilation algorithm}\label{sec:compile}

We are now ready to state the algorithm. The input is a finite GSLT
presentation $\Pi$; the output is a Rholang program $\mathsf{compile}(\Pi)$.
The compiler is presented as nine emission rules, each of which produces a
piece of Rholang code parameterised over data extracted from $\Pi$ and the
set automaton $M = (\States, s_0, L, \delta, \eta)$ built from $\Lhs(\Pi)$.

\begin{algorithm}[$\mathsf{compile}(\Pi)$]\label{alg:compile}
\hfill\\
\noindent
\textbf{Input:} a finite GSLT presentation
$\Pi = (\mathcal{S}, \mathbb{F}, E, \succ, \mathcal{R}_{\mathrm{rule}}, \mathcal{R}_{\mathrm{ctx}})$.\\
\textbf{Output:} a closed Rholang program.\\
\noindent
\textbf{Steps.}
\begin{enumerate}[leftmargin=2em]
  \item \emph{Normalise equations.} Use Knuth--Bendix on $(E, \succ)$ to
        obtain a confluent, terminating rewrite system $E^{\downarrow}$
        whose normal forms are the canonical representatives modulo $E$.
        The remaining steps assume that all left-hand sides in
        $\mathcal{R}_{\mathrm{rule}} \cup \mathcal{R}_{\mathrm{ctx}}$ are
        in $E^{\downarrow}$-normal form.
  \item \emph{Build $\Lhs$.} Collect the left-hand sides of
        $\mathcal{R}_{\mathrm{rule}}$ and the per-premise left-hand sides
        of $\mathcal{R}_{\mathrm{ctx}}$.
  \item \emph{Build the set automaton $M$.} Compute the states, $L$,
        $\delta$, and $\eta$ by the construction of
        \S\ref{ssec:set-automaton}.
  \item \emph{Emit the runtime preamble} (\S\ref{ssec:emit-preamble}).
  \item \emph{Emit the term constructors} (\S\ref{ssec:emit-constructors}).
  \item \emph{Emit the traversal contract} (\S\ref{ssec:emit-traversal}).
  \item \emph{Emit the state contracts} (\S\ref{ssec:emit-states}).
  \item \emph{Emit the equational normaliser} (\S\ref{ssec:emit-normaliser}).
  \item \emph{Emit the rewriter} (\S\ref{ssec:rewriter}).
  \item \emph{Emit the driver} (\S\ref{ssec:driver}).
\end{enumerate}
The order of emission is not significant; emitted contracts are persistent
and unordered.
\end{algorithm}

We describe each emission rule in turn. All emitted code lives inside a
single top-level \texttt{new} block that introduces the channel names used
throughout. We elide that wrapper in displayed code.

\subsection{Preamble}\label{ssec:emit-preamble}

The preamble declares the public channels:
\begin{itemize}[leftmargin=*]
  \item \texttt{matches} --- a public channel on which the interpreter
        announces every $(L, p)$-pair such that $L$ matches the current
        subject term at position $p$;
  \item \texttt{fire} --- a public channel on which a scheduler may
        instruct the interpreter to fire a specific announced redex;
  \item \texttt{normalForm} --- a public channel that receives the final
        term when no more redexes are available.
\end{itemize}
\begin{lstlisting}
new matches, fire, normalForm,
    /* state channels: */ ${state-names},
    /* helpers: */ traverse, substitute, normalise, driver in {
  // ... remainder of the program ...
}
\end{lstlisting}
The placeholder \texttt{\$\{state-names\}} expands to the list
$\sigma_{s_0}, \sigma_{s_1}, \ldots, \sigma_{s_{|\States|-1}}$.

\subsection{Term constructors}\label{ssec:emit-constructors}

For each $f \in \mathbb{F}_n$ the compiler emits a \emph{constructor
contract} that builds a term-channel responder for $f$:
\begin{lstlisting}
for (@(t1, ..., tn), r <= mk_${f}) {
  new tau in {
    for (reply <= tau) { reply!((${f-tag}, t1, ..., tn)) } |
    r!(tau)
  }
}
\end{lstlisting}
The placeholder \texttt{\$\{f-tag\}} is the canonical name $\reflect{f}$.
A constant $a \in \mathbb{F}_0$ emits the obvious zero-argument variant.

For each $f \in \mathbb{F}_n$ with binding positions $\beta_f$, the
$i \in \beta_f$ slot is a name $\tau_i$ that satisfies the bound-body
protocol described in \S\ref{ssec:term-encoding}; the constructor contract
above accepts such a name without further change.

\paragraph{Well-formedness equations.}
When $\lambda_f$ is non-empty, the constructor contract guards
$f$-formation on the input equations: it invokes
\texttt{checkConstraint} (\S\ref{ssec:rewriter}) on the inputs
$(t_1, \ldots, t_n)$ against $\lambda_f$ before publishing the new
term channel, and emits a typing failure on $r$ if any equation fails.
The check is structurally identical to the sharing-constraint check
of \S\ref{ssec:rewriter}: each equation $g_i(\vec{x}) = h_i(\vec{x})$
is evaluated to two term channels and compared by Rholang structural
congruence. When $\lambda_f = \emptyset$ the guard is elided and the
contract reduces to the form above.

\subsection{Traversal contract}\label{ssec:emit-traversal}

A single \emph{traversal contract} consumes a term channel and a position
list and replies with the term channel at that position:
\begin{lstlisting}
for (@(tau, p), reply <= traverse) {
  match p {
    Nil       => { reply!(tau) }
    [i, ...q] => {
      new r in {
        tau!(r) |
        for ((_head, @args) <- r) { traverse!((args.nth(i), q), reply) }
      }
    }
  }
}
\end{lstlisting}
\emph{Note.} The pattern \lstinline|(_head, @args)| extracts the head tag
and the tuple of argument channels; \lstinline|args.nth(i)| selects the
$i$-th of them. The traversal is asynchronous: every step requests the
next argument channel via a fresh response name $r$ and resumes only
when the term channel replies.

\subsection{State contracts}\label{ssec:emit-states}

The set-automaton dispatch is the heart of the compiler. For each
non-final state $s \in \States$ we emit a single persistent contract on
the state channel $\sigma_s$:
\begin{lstlisting}
for (@(tau, p) <= ${state-name(s)}) {
  new r, posChan in {
    // 1. Compute the inspection position p . L(s):
    posChan!(p ++ ${pos-literal(L(s))}) |
    for (@inspectPos <- posChan) {
      // 2. Ask traverse for the term channel at inspectPos:
      traverse!((tau, inspectPos), r) |
      for (subTau <- r) {
        new rr in {
          subTau!(rr) |
          for ((@head, @args) <- rr) {
            // 3. Dispatch on the head symbol:
            match head {
              ${case-for-each-f-in-F(s)}
              _ => Nil
            }
          }
        }
      }
    }
  }
}
\end{lstlisting}
The placeholder \texttt{\$\{case-for-each-f-in-F(s)\}} expands to one
\texttt{match} arm per function symbol $f$ for which $\delta(s, f) \cup
\eta(s, f)$ is non-empty. The arm has the form
\begin{lstlisting}
  ${f-tag} => {
    // (a) emit pattern-match announcements eta(s, f):
    ${emit-eta(s, f)} |
    // (b) emit transitions delta(s, f), in parallel:
    ${emit-delta(s, f)}
  }
\end{lstlisting}
where
\[
  \mathtt{emit\text{-}eta}(s, f) \;=\;
  \bigm|_{(\ell, p', \kappa) \in \eta(s, f)}
  \mathtt{matches}\,!\bigl(\,\reflect{\ell},\; p \mathbin{+\!+} p',\;
                            \reflect{\kappa\langle p.p'\rangle}\,\bigr)
\]
and
\[
  \mathtt{emit\text{-}delta}(s, f) \;=\;
  \bigm|_{(s', p'') \in \delta(s, f)}
  \sigma_{s'}\,!(\,\tau,\; p \mathbin{+\!+} p''\,).
\]
The notation $\kappa\langle q \rangle$ denotes $\kappa$ with every
position-pair $(q_1, q_2)$ replaced by $(q.q_1, q.q_2)$ (the
reanchoring at the announcement position). When $\kappa$ is empty
(i.e., the LHS is syntactically linear), $\reflect{\kappa\langle \cdot \rangle}$
is a canonical name for the empty constraint, and the fire contract
(\S\ref{ssec:rewriter}) bypasses the equality check.
The vertical bar denotes Rholang parallel composition: distinct
announcements and distinct successor configurations are emitted in
parallel and are scheduled by the Rholang runtime.

The final state (the empty set of match goals) emits no contract: a
configuration that arrives at the final state simply terminates that
branch of the search.

\paragraph{What this code is doing.}
Each state contract reads a configuration $(\tau, p)$ and:
\begin{enumerate}[leftmargin=*]
  \item computes the inspection position $p.L(s)$;
  \item asks the traversal contract for the term channel at that
        position;
  \item asks that channel for its head symbol $f$ and argument
        channels;
  \item announces all matches $\eta(s, f)$ in parallel;
  \item dispatches all successor configurations $\delta(s, f)$ in
        parallel.
\end{enumerate}
There is no shared state and no scheduling barrier. The arrows of
$\delta(s, f)$ correspond to distinct messages on distinct or identical
channels, all available simultaneously to any Rholang reducer.

\subsection{Equational normaliser}\label{ssec:emit-normaliser}

The equations $E$ of the presentation, after Knuth--Bendix completion,
form a convergent rewrite system $E^{\downarrow}$ that gives a canonical
representative for each equivalence class. Because Definition~\ref{def:presentation}
requires $E^{\downarrow}$ to satisfy the finite variant property, a
subject term and its $E$-canonical form admit the \emph{same} set of
pattern matches against any left-hand side of $\mathcal{R}$ (in the
sense of $E$-matching): normalising before dispatch is sound and
complete. The compiler treats $E^{\downarrow}$ as a sub-presentation
whose set automaton $M^{\downarrow}$ matches the left-hand sides of
$E^{\downarrow}$ in the subject term, and whose rewriter
(\S\ref{ssec:rewriter}) fires them eagerly. We emit a \emph{normaliser}
contract that, given a term channel, repeatedly fires
$E^{\downarrow}$-redexes until none remain and then signals completion:
\begin{lstlisting}
for (tau, done <= normalise) {
  // Run the E-downward set automaton, fire its redexes innermost-first,
  // and signal "done" when matches is exhausted on E-downward patterns.
  // ...
}
\end{lstlisting}
The full body is structurally identical to the driver in
\S\ref{ssec:driver}; we omit it. Normalisation is interleaved with
$\mathcal{R}$-rewriting in the driver.

\subsection{Rewriter}\label{ssec:rewriter}

For each rule $L_j \rew R_j \in \mathcal{R}_{\mathrm{rule}}$ the compiler
emits a \emph{fire contract} parameterised by the pattern identifier
$\reflect{L_j}$. The contract takes the announcement triple
$(\tau, p, \kappa)$ produced by the matcher; if $\kappa$ is non-empty,
the contract first verifies the structural equalities it specifies and
aborts the firing if any fails. The Rholang skeleton is:
\begin{lstlisting}
for (@(tau, p, kappa), done <= fire_${L_j}) {
  new sigma, rTau, eqOk in {
    // 0. Verify sharing constraints kappa, if any:
    checkConstraint!((tau, kappa), eqOk) |
    for (@ok <- eqOk) {
      if (ok) {
        // 1. Extract the substitution sigma from tau at p against L_j:
        extractSubst!((tau, p, ${pattern-id(L_j)}), sigma) |
        for (@s <- sigma) {
          // 2. Build the right-hand side R_j under sigma on a fresh channel:
          buildRhs!((${rhs-id(R_j)}, @s), rTau) |
          // 3. Splice rTau into tau at p:
          for (newTau <- rTau) {
            splice!((tau, p, newTau), done)
          }
        }
      } else {
        done!(@"abort")
      }
    }
  }
}
\end{lstlisting}
The auxiliary contracts \texttt{checkConstraint}, \texttt{extractSubst},
\texttt{buildRhs}, and \texttt{splice} are generic and emitted once each.

\paragraph{Constraint check.}
The \texttt{checkConstraint} contract is the workhorse for every
equality side-condition emitted by the compiler. It receives a
term channel $\tau$ and a finite set of \emph{equality items}
$\xi = \{(e_1, e'_1), \ldots, (e_m, e'_m)\}$ and replies \texttt{true}
iff $e_i \equiv e'_i$ in $\tau$ for every $i$, where each $e_i$ is
either:
\begin{itemize}[leftmargin=*]
  \item a \emph{position pair}, as used in sharing constraints $\kappa$
        and resolved by \texttt{traverse} into a sub-channel of $\tau$; or
  \item an $\mathbb{F}$-term over the inputs of a constructor, as used in
        well-formedness equations $\lambda_f$ and resolved by
        \texttt{buildRhs} into a freshly built term channel.
\end{itemize}
In both cases, each pair-check is realised by two parallel resolution
steps returning channels $c_i, c'_i$ and a single Rholang \texttt{match}
that succeeds iff $c_i \equiv c'_i$ as quoted processes (a decidable
test on rho-calculus structural congruence). The $m$ pair-checks run
in parallel; \texttt{checkConstraint} joins their booleans by
conjunction. The constructor contracts of \S\ref{ssec:emit-constructors}
invoke \texttt{checkConstraint} with $\xi = \lambda_f$ resolved on the
inputs; the fire contracts above invoke it with $\xi = \kappa$
resolved on the subject channel.

\paragraph{Substitution and binders.}
The auxiliary \texttt{buildRhs} contract takes a substitution $\sigma$
(a mapping from metavariables of $L_j$ to term channels) and the static
syntactic data of $R_j$, and produces a term channel for $R_j^{\sigma}$.
For a binding occurrence in $R_j$ --- a metavariable $\omega$ that occurs
under a binding position of some constructor --- the corresponding
channel from $\sigma$ already carries a bound-name token in its
protocol; \texttt{buildRhs} composes it with the surrounding constructor
contracts straightforwardly.

\paragraph{Contextual rewrites.}
A contextual rewrite
\[
  L_1 \rew R_1, \ldots, L_k \rew R_k \Rightarrow K[L_1, \ldots, L_k] \rew K'[R_1, \ldots, R_k]
\]
is compiled in two stages. The premise patterns $L_i$ are included in
$\Lhs$ and matched by the set automaton as in any other rule. The
contextual conclusion fires only when announcements of all $k$ premises
arrive on \texttt{matches} in a configuration that places them at the
$k$ holes of $K$; this synchronisation is implemented as a \emph{join} on
$k$ messages on the \texttt{matches} channel, gated by a check that the
positions are compatible with $K$. The check is itself a finite-state
predicate on positions, emitted by the compiler from $K$.

\subsection{Driver}\label{ssec:driver}

The driver wires the components together. Given an initial subject term
$t_0$, the driver:
\begin{enumerate}[leftmargin=*]
  \item builds the term channel $\tau_0$ for $t_0$ via the constructor
        contracts;
  \item normalises $\tau_0$ with the equational normaliser;
  \item kicks off the set automaton by sending $(\tau_0, \epsilon)$ on
        $\sigma_{s_0}$;
  \item collects announcements on \texttt{matches} into a redex pool;
  \item upon receiving an instruction on \texttt{fire}, dispatches the
        appropriate fire contract;
  \item after each firing, restarts the automaton from $\sigma_{s_0}$
        on the new term channel (or, with a more refined construction,
        from the partial-evaluation state determined by the rewritten
        position);
  \item when \texttt{matches} reports empty after a full automaton run,
        sends the current term on \texttt{normalForm} and terminates.
\end{enumerate}
The default scheduling policy is to forward every announcement back as
an instruction to fire, in some order; alternative policies (outermost,
innermost, etc.) are obtained by filtering or sorting the redex pool
before forwarding.

\section{Parallelism preservation}\label{sec:parallel}

The defining feature of the set automaton is that it exposes parallelism
in two places: \emph{intra-state} (the elements of $\delta(s, f)$ and
$\eta(s, f)$ are independent and may be processed in any order) and
\emph{inter-rule} (announced matches at non-overlapping positions are
independent redexes). We show that both kinds of parallelism are
preserved verbatim by Algorithm~\ref{alg:compile}.

\subsection{Concurrency model}

Recall that the Rholang runtime is asynchronous and name-passing: a send
$\Send{x}{P}$ may be paired with any matching receive on $x$, in any
order, and parallel composition $P_1 \mid P_2$ enforces no ordering. The
reduction relation of the rho calculus is closed under parallel
composition: $P \rew P'$ implies $P \mid Q \rew P' \mid Q$. Two parallel
sub-processes may reduce independently and concurrently, so long as
their reductions do not compete for the same linear receive.

\begin{definition}[Independent reductions]\label{def:independent}
Two reductions $P_1 \rew P_1'$ and $P_2 \rew P_2'$ of sub-processes of
$P$ are \emph{independent} if they involve disjoint sets of names and
disjoint sub-processes.
\end{definition}

\subsection{Intra-state parallelism}

\begin{theorem}[Intra-state parallelism]\label{thm:intra}
Let $s$ be a state of the set automaton and $f$ a function symbol such
that $\delta(s, f) = \{(s'_1, p''_1), \ldots, (s'_m, p''_m)\}$ and
$\eta(s, f) = \{(\ell_1, q_1), \ldots, (\ell_k, q_k)\}$. After dispatch
on the head symbol $f$ at a configuration $(\tau, p)$, the corresponding
Rholang sub-process is
\[
  \bigm|_{j=1}^{k} \texttt{matches!}((\reflect{\ell_j}, p.q_j))
  \;\mathbin{\Big|}\;
  \bigm|_{i=1}^{m} \texttt{$\sigma_{s'_i}$!}((\tau, p.p''_i)).
\]
The $k + m$ resulting message sends are pairwise independent in the sense
of Definition~\ref{def:independent}; the Rholang runtime may schedule
them in any order or concurrently, and the resulting global behaviour is
identical to that of the set automaton run with any sequential schedule
of the corresponding $(s', p)$-successors and $(\ell, p)$-announcements.
\end{theorem}

\begin{proof}
The $k + m$ message sends share only the read-only names $\tau$ and $p$
(the latter as a fresh local list, not a Rholang channel target); the
target channels \texttt{matches} and $\sigma_{s'_1}, \ldots, \sigma_{s'_m}$
are pairwise distinct or, in the case that two transitions go to the
same state $s'_i = s'_j$, both receive messages on the same persistent
contract --- which by definition processes them in some order without
interference. Soundness of any scheduling order is exactly the
soundness theorem for the set automaton~\cite[Theorem~6.1]{ErkensGroote2021};
that the Rholang runtime may schedule any order is immediate from the
asynchronous semantics of $\mid$ and the persistence of the state
contracts.
\end{proof}

\subsection{Inter-rule parallelism}

\begin{theorem}[Inter-rule parallelism]\label{thm:inter}
Let $\Pi$ be a finite GSLT presentation and let $t$ be a subject term.
Let
\[
  A \;=\; \{(\ell, p) : \text{\textnormal{the set automaton announces $\ell$ at $p$ during the run on $t$}}\}
\]
be the set of announcements produced by the compiled Rholang program in
its first \emph{automaton phase}. The Rholang process announcing $A$ is
of the form
\[
  \bigm|_{(\ell, p) \in A} \texttt{matches!}((\reflect{\ell}, p)) \;\mid\; (\text{automaton continuations}),
\]
and the announcements are pairwise independent. Any subset of redexes
fired by the driver in parallel (a \emph{parallel firing}
$\{(\ell_{i_1}, p_{i_1}), \ldots, (\ell_{i_r}, p_{i_r})\}$ with pairwise
non-overlapping positions) is realised by the Rholang runtime as the
parallel firing of $r$ independent fire-contract invocations, without
inter-firing scheduling glue.
\end{theorem}

\begin{proof}
The set of announcements is the parallel composition of independent
sends to \texttt{matches} by Theorem~\ref{thm:intra} composed
transitively over the automaton run. The driver dispatches each fire
instruction to a fresh \texttt{fire\_$\ell$} contract invocation, which
allocates its own \texttt{rTau}, \texttt{sigma}, and other local
channels via \texttt{new}, so no two fire contracts share linear
channels. Non-overlapping positions imply non-overlapping subject-term
sub-channels, which by Construction~\ref{cons:term-encoding} are
allocated under distinct \texttt{new}-restrictions. Hence the $r$
firings touch disjoint sets of names and disjoint sub-processes, and
proceed independently.
\end{proof}

\subsection{Why this matters}

The combined effect of Theorems~\ref{thm:intra} and~\ref{thm:inter} is
that no part of the compiled interpreter contains a centralised
scheduler. The matcher emits successor configurations as parallel sends;
the rewriter emits parallel firings as parallel sub-processes; the
driver merely forwards. Concretely, the only points at which the
interpreter is sequential are the points at which the set automaton is
inherently sequential --- namely, the inspection of the head symbol of a
single sub-term at a single position. Every other step is parallel by
construction, and Rholang's runtime is responsible for choosing any
interleaving.

\begin{corollary}[No-glue corollary]\label{cor:no-glue}
The compiled Rholang program contains no parallel-scheduling code: the
emitted runtime preamble, traversal contract, state contracts, fire
contracts, and driver are all sequential within their bodies, and all
inter-component parallelism is exposed solely through Rholang's
$\mid$-composition. In particular, the size of the compiled program is
$O(|\States| \cdot |\mathbb{F}| + |\mathcal{R}| + |\mathbb{F}|)$, with no
multiplicative factor for scheduling.
\end{corollary}

\section{A worked example}\label{sec:example}

We illustrate the algorithm on a tiny presentation: the associativity
patterns of Erkens--Groote~\cite[\S3]{ErkensGroote2021}, treated as a
rewrite system that re-associates a binary operator to the right.

\subsection{The presentation}

Sorts: $\mathcal{S} = \{\iota\}$. Signature: $\mathbb{F}_0 = \{a\}$,
$\mathbb{F}_2 = \{f\}$. Equations: none beyond the trivial. Rules:
\[
  \mathcal{R}_{\mathrm{rule}} = \{\; L_1 \rew R_1 \;\}
  \quad\text{where}\quad
  L_1 = f(f(\omega_1, \omega_2), \omega_3),\;
  R_1 = f(\omega_1, f(\omega_2, \omega_3)).
\]
Contextual rewrites: none. The left-hand-side set is
$\Lhs = \{L_1\}$. (The Erkens--Groote example uses a second pattern
$L_2 = f(\omega_1, f(\omega_2, \omega_3))$ to demonstrate set automata;
we drop it here so the worked example fits on one page.)

\subsection{The set automaton}

Applying the construction of \S\ref{ssec:set-automaton} to
$\Lhs = \{L_1\}$ yields:
\begin{itemize}[leftmargin=*]
  \item Initial state $s_0 = \{L_1 @ \epsilon \to L_1 @ \epsilon\}$ with
        $L(s_0) = \epsilon$;
  \item $f$-derivative of $s_0$ contains the new goal $f(\omega_1, \omega_2) @ 1 \to L_1 @ \epsilon$
        and fresh goals $L_1 @ 1 \to L_1 @ 1$, $L_1 @ 2 \to L_1 @ 2$;
        partitioning and lifting yield successors at position pointer
        $1$ and $2$;
  \item further derivatives yield a finite set of states with
        position labels in $\{\epsilon, 1, 2\}$ and announcements
        $\eta(s, f) = \{(L_1, \epsilon)\}$ at the (unique) state $s$
        whose obligations require both $\hd@\epsilon = f$ and
        $\hd@1 = f$.
\end{itemize}
We abbreviate the resulting four-state automaton as
$\States = \{s_0, s_1, s_2, s_3\}$.

\subsection{The compiled Rholang program}

We give the key fragments. The preamble:
\begin{lstlisting}
new matches, fire, normalForm,
    s0, s1, s2, s3,           // state channels
    mk_a, mk_f,                // constructors
    traverse, fire_L1,
    extractSubst, buildRhs, splice,
    driver in {
  // ... contracts below ...
  driver!()
}
\end{lstlisting}
The constructors:
\begin{lstlisting}
for (r <= mk_a) {
  new tau in {
    for (reply <= tau) { reply!((@"a",)) } |
    r!(tau)
  }
} |
for (@(t1, t2), r <= mk_f) {
  new tau in {
    for (reply <= tau) { reply!((@"f", t1, t2)) } |
    r!(tau)
  }
}
\end{lstlisting}
The state contract for $s_0$ (label $\epsilon$, $f$-transition to $s_1$
at position $1$ and $s_2$ at position $2$, with the announcement
$(L_1, \epsilon)$ deferred to a deeper state):
\begin{lstlisting}
for (@(tau, p) <= s0) {
  new r in {
    traverse!((tau, p), r) |        // L(s_0) = epsilon, so p ++ [] = p
    for (subTau <- r) {
      new rr in {
        subTau!(rr) |
        for ((@head, @args) <- rr) {
          match head {
            @"f" => {
              s1!((tau, p ++ [1])) |  // descend to position p.1
              s2!((tau, p ++ [2]))    // descend to position p.2
            }
            _    => Nil
          }
        }
      }
    }
  }
}
\end{lstlisting}
The state contract for the state $s$ at which the announcement fires:
\begin{lstlisting}
for (@(tau, p) <= s1) {
  new r in {
    traverse!((tau, p), r) |        // L(s_1) = epsilon
    for (subTau <- r) {
      new rr in {
        subTau!(rr) |
        for ((@head, @args) <- rr) {
          match head {
            @"f" => {
              // announce: L_1 matches at p.epsilon (= p)
              matches!((@"L_1", p)) |
              // and continue the inner search:
              s1!((tau, p ++ [1])) |
              s2!((tau, p ++ [2]))
            }
            _    => Nil
          }
        }
      }
    }
  }
}
\end{lstlisting}
(\emph{Note.} The redex announcement and the inner search are emitted in
parallel. The Rholang runtime is free to deliver them in any order; both
will eventually run.)

The fire contract for $L_1$, sketched:
\begin{lstlisting}
for (@(tau, p), done <= fire_L1) {
  new sigma, rTau in {
    extractSubst!((tau, p, @"L_1"), sigma) |
    for (@s <- sigma) {
      // s = { omega_1 |-> t11, omega_2 |-> t12, omega_3 |-> t2 } where
      // t = f(f(t11, t12), t2)
      mk_f((s.omega_2, s.omega_3), @"inner") |
      for (innerTau <- @"inner") {
        mk_f((s.omega_1, innerTau), rTau) |
        for (newTau <- rTau) {
          splice!((tau, p, newTau), done)
        }
      }
    }
  }
}
\end{lstlisting}
(In a production compiler the literal name \texttt{@"inner"} would be a
fresh \texttt{new}-bound name; we use the literal here for readability.)

The driver simply forwards every announcement to a fire instruction and
restarts the automaton:
\begin{lstlisting}
for (_ <= driver) {
  // ... build initial tau_0 from input ...
  s0!((tau_0, [])) |
  for (@(L, p) <= matches) {           // persistent for: process every announcement
    fire_L1!((tau_0, p), @"fired") |
    for (_ <- @"fired") { s0!((tau_0, [])) }
  }
}
\end{lstlisting}

\subsection{Running it}

On the subject term $t = f(f(a, a), a)$ the program proceeds as
follows.
\begin{itemize}[leftmargin=*]
  \item The driver builds the term channel $\tau_0$ for $t$ and sends
        $(\tau_0, [\,])$ on $\sigma_{s_0}$.
  \item $\sigma_{s_0}$ inspects the head at position $\epsilon$,
        finds $f$, and emits two parallel sends:
        $\sigma_{s_1}(\tau_0, [1])$ and $\sigma_{s_2}(\tau_0, [2])$.
  \item $\sigma_{s_1}(\tau_0, [1])$ inspects the head at position $1$,
        finds $f$, announces $(L_1, [\,])$, and emits two further
        parallel sends.
  \item $\sigma_{s_2}(\tau_0, [2])$ inspects the head at position $2$,
        finds $a$, and terminates (no $a$-arm in any state contract).
  \item The announcement $(L_1, [\,])$ is forwarded by the driver to
        \texttt{fire\_L1}, which extracts $\sigma$, builds the RHS,
        splices it into $\tau_0$, and restarts the automaton.
\end{itemize}
Each function symbol of $t$ is inspected exactly once per automaton run,
and the two parallel sends from $\sigma_{s_0}$ are scheduled by the
Rholang runtime in either order or simultaneously.

\subsection{The rho COMM rule: full compilation pipeline}\label{ssec:rho-comm}

We now revisit the rho-calculus COMM rule of \S\ref{ssec:set-automaton},
filling in the compilation pipeline downstream of the set automaton. The
linearisation
$\ell^{\flat}_{\mathrm{COMM}}$ and the sharing constraint
$\kappa_{\mathrm{COMM}} = \{(1.1,\, 2.1)\}$ were derived there; we take
them as given here and trace what happens between the matcher's
announcement and the firing of the redex.

At announcement, the state-contract template of \S\ref{ssec:emit-states}
emits
\[
  \mathtt{matches}\,!\bigl(\,\reflect{\mathrm{COMM}},\; p,\;
                             \reflect{\{(p.1.1,\, p.2.1)\}}\,\bigr).
\]
The driver (\S\ref{ssec:driver}) forwards the triple to the fire
contract \texttt{fire\_COMM}, which:
\begin{enumerate}[leftmargin=*]
  \item invokes \texttt{checkConstraint} on the pair $(p.1.1,\, p.2.1)$,
        which fires two parallel \texttt{traverse} calls returning
        sub-channels $c_1, c_2$ of the subject term, then asks the
        Rholang runtime whether $c_1 \equiv c_2$ as quoted processes;
  \item if the check fails, aborts the firing and yields control back
        to the driver, which moves on to the next announcement;
  \item if the check succeeds, calls \texttt{extractSubst} on
        $(\tau, p, \reflect{\mathrm{COMM}})$ to extract a substitution
        $\sigma$ mapping the metavariables of the original
        (unlinearised) LHS to subject-term channels --- here, $x$ to
        $c_1$ (equivalently $c_2$), $B$ to the body channel at $p.1.2$,
        and $Q$ to the payload channel at $p.2.2$;
  \item calls \texttt{buildRhs} on the static RHS data of the COMM rule
        and $\sigma$, which produces a fresh term channel for the
        process $B[@Q/y]$ by composing the constructor contracts of the
        rho-calculus signature; the binder $y$ in the RHS is handled by
        the bound-body protocol of \S\ref{ssec:term-encoding};
  \item calls \texttt{splice} on $(\tau, p, \text{new channel})$, which
        replaces the sub-channel of $\tau$ at $p$ by the new channel and
        signals completion on \texttt{done}.
\end{enumerate}
The driver receives the completion signal and restarts the automaton on
the updated subject term (or, with a more refined construction, restarts
from a partial-evaluation state determined by the rewritten position).
The set automaton itself never inspects whether the two channels are
equal; that information lives entirely in the fire stage, where it is
decidable by the host structural-congruence check. Likewise, the host
runtime is responsible for choosing an interleaving of the parallel
\texttt{traverse}-calls inside \texttt{checkConstraint}, the parallel
sends to successor states, and the parallel fire-contract invocations
on concurrent announcements.

\section{Discussion}\label{sec:discussion}

\subsection{Garbage collection}

Term channels and state-channel messages are reference-counted naturally
by Rholang's structural congruence: a name with no pending sends and no
pending receivers is structurally equivalent to $\nil$ and is collected
by the runtime. Our encoding does not introduce reference cycles
(constructor contracts are leaves; state contracts are persistent but
fixed in number; fire contracts are linear and consume themselves), so
the standard rho-calculus garbage collector applies.

\subsection{Dynamic rule sets}

The compiler as stated treats the rule set as static. Adding rules at
runtime --- adding new state-contract code --- is straightforward in
Rholang: the new state contracts can be emitted as additional
\texttt{contract} declarations on existing channels. The set automaton,
however, must be rebuilt: adding a new pattern can change the dependency
partition at every reachable state. A practical implementation would
recompute the affected states only, but a full discussion is beyond
this paper.

\subsection{Equations, confluence, and FVP}

The compiler's correctness rests on two assumptions about $E$:
convergence of $E^{\downarrow}$ (so that canonical forms exist and are
unique) and the finite variant property (so that $E$-matching against
the LHSs of $\mathcal{R}$ is decidable and reduces to syntactic
matching after normalisation). When convergence fails --- $E$ is
non-orientable, or KB completion does not terminate --- the compiler
emits a warning and falls back to an interleaved $E$-and-$\mathcal{R}$
rewriting strategy that is sound on confluent inputs but offers no
completeness guarantee. When convergence holds but FVP fails --- for
example, equations like $f(f(x)) = f(x)$ in unrestricted form, which
generate infinitely many variants --- $E$-matching is undecidable in
general, and the compiler can no longer guarantee that the set
automaton, built on $E$-canonical LHSs, finds every $E$-equivalent
redex. Sound but incomplete behaviour can be recovered by restricting
to syntactic matching against canonical forms; complete behaviour
requires either FVP, one of the other syntactic conditions listed in
the footnote to Definition~\ref{def:presentation}, or a manually
provided $E$-matching oracle.

\subsection{Comparison to direct compilation}

A naive compiler from a GSLT presentation to Rholang would emit, for
each rule $L_j \rew R_j$, a single \texttt{for} that pattern-matches
$L_j$ on the entire subject term. Such a compiler trades program size
for redundant inspection: each function symbol of the subject term is
inspected once per rule per ancestor position. The set-automaton route
amortises this cost across the rule set, inspecting each function
symbol once \emph{total} per automaton run, at the price of a richer
state space. For pattern sets of even modest size the saving is
substantial, and the parallelism gained is, as we have shown, free.

\section{Conclusion}\label{sec:conclusion}

The algorithm $\mathsf{compile}$ packages three independent strands of
work --- finitely presented graph-structured lambda theories, the
Erkens--Groote set automaton, and the rho calculus --- into a single
compilation pipeline. The output is a Rholang interpreter that exposes
every parallelism opportunity of the set automaton through Rholang's
asynchronous, name-passing semantics, without scheduling glue. We have
shown the compiler is total, polynomial in the size of the automaton,
and parallelism-preserving by construction. The pipeline gives a
concrete answer to the question of how to implement, in a concurrent
process calculus, the small-step semantics of any finitely presented
GSLT.

\paragraph{Future work.}
Several directions remain. Incremental rebuild of the set automaton
under dynamic rule changes would make the pipeline suitable for
language-development environments. Partial evaluation of the automaton
against statically known sub-contexts of the subject term would shrink
the automaton --- and the compiled program --- in common cases such as
pre-compiled standard libraries. Finally, a formal mechanisation of the
parallelism-preservation theorem in a proof assistant would lend the
construction the same kind of foundational credibility that the
set-automaton construction enjoys in~\cite{ErkensGroote2021}.

\begin{thebibliography}{9}

\bibitem{ErkensGroote2021}
R.~Erkens and J.~F.~Groote.
\newblock {A set automaton to locate all pattern matches in a term}.
\newblock \emph{arXiv preprint 2106.15311}, 2021.

\bibitem{MeredithRadestock2005}
L.~G.~Meredith and M.~Radestock.
\newblock {A reflective higher-order calculus}.
\newblock In \emph{Electronic Notes in Theoretical Computer Science},
  volume 141(5), pages 49--67, 2005.

\bibitem{EscobarMeseguerSasse2008}
S.~Escobar, J.~Meseguer, and R.~Sasse.
\newblock {Effectively checking the finite variant property}.
\newblock In \emph{Rewriting Techniques and Applications (RTA)},
  LNCS 5117, pages 79--93, 2008.

\bibitem{EscobarSasseMeseguer2012}
S.~Escobar, R.~Sasse, and J.~Meseguer.
\newblock {Folding variant narrowing and optimal variant termination}.
\newblock \emph{Journal of Logic and Algebraic Programming}, 81(7--8),
  pages 898--928, 2012.

\bibitem{TrimbleNLabLawvere}
T.~Trimble.
\newblock {Multisorted Lawvere theories}.
\newblock nLab entry, revision of 2024.
\newblock \url{https://ncatlab.org/toddtrimble/published/multisorted+Lawvere+theories}.

\end{thebibliography}

\end{document}
